\documentclass{article}
\usepackage{lmodern}
\usepackage{amsmath}
\usepackage{listings}
\usepackage{fullpage}
\usepackage{parskip}
\usepackage{xcolor}
\lstset{
	basicstyle=\ttfamily,
	columns=fullflexible,
	frame=single,
	breaklines=true,
	postbreak=\mbox{\textcolor{red}{$\hookrightarrow$}\space},
}
\begin{document}
\title{Experiment `p\_6\_10\_same\_set\_array' Results}
\author{\today}
\date{}
\maketitle
\textbf{Experiment outcome: }SUCCESS
\\\textbf{Bad responses: }0
\\\textbf{Responses containing}~\texttt{assume}~\textbf{: }0
\\\textbf{Responses containing}~\texttt{expect}~\textbf{: }0
\\\textbf{Resolution attempts: }1
\\\textbf{Hard fails (resolution): }0
\\\textbf{Soft fails (resolution): }0
\\\textbf{Verification attempts: }1
\section*{Problem Specification}
\textbf{Problem name: }p\_6\_10\_same\_set\_array
\\\textbf{Natural language statement: }Write a method that checks whether two arrays have the same elements in some order, ignoring duplicates.
\\\textbf{Method signature: }p\_6\_10\_same\_set\_array(first: array<int>, second: array<int>) returns (are\_same\_set: bool)
\subsection*{Ensures}
\begin{itemize}
\item \texttt{are\_same\_set == (array\_to\_set(first[..]) == array\_to\_set(second[..]))}
\end{itemize}
\subsection*{Functional Code Given}
\begin{lstlisting}
function array_to_set(arr: seq<int>): set<int>
{
  if |arr| == 0 then {}
  else {arr[0]} + array_to_set(arr[1..])
}
\end{lstlisting}
\clearpage
\section*{GenAI interactions}
Below you will find all interactions between the `user' (program) and the `assistant' (OpenAI).
\subsection*{Program $\rightarrow$ GenAI}
\begin{lstlisting}
You are given the following task to perform in Dafny:

Write a method that checks whether two arrays have the same elements in some order, ignoring duplicates.

The signature should be:

p_6_10_same_set_array(first: array<int>, second: array<int>) returns (are_same_set: bool)

The method should respect the following contract:

ensures are_same_set == (array_to_set(first[..]) == array_to_set(second[..]))

The contract uses the following dafny code:

function array_to_set(arr: seq<int>): set<int>
{
  if |arr| == 0 then {}
  else {arr[0]} + array_to_set(arr[1..])
}

These function(s) must not be used in your implementation of the method.

Produce and show only the Dafny body of this method, including the curly braces that surround it. Do not show the signature nor contract. You must not use 'assume' nor 'expect' anywhere in your code.
\end{lstlisting}
\subsection*{GenAI $\rightarrow$ Program}
\textbf{System fingerprint: }null
\\\textbf{ID: }chatcmpl-Dn2STVyMhUNtLoMw5s4CUMeBYxtW4
\\\textbf{Model: }gpt-5.5-2026-04-23
\\\textbf{Created at: }1780579893
\\\textbf{Finish reason: }stop
\\\textbf{Usage: }434 tokens in, and 6186 tokens out
\begin{lstlisting}
{
  var firstSet: set<int> := {};
  var i := first.Length;
  while i > 0
    invariant 0 <= i <= first.Length
    invariant firstSet == array_to_set(first[i..])
    decreases i
  {
    var k := i - 1;
    assert 0 <= k < first.Length;
    assert firstSet == array_to_set(first[k + 1..]);
    assert |first[k..]| > 0;
    assert first[k..][0] == first[k];
    assert first[k..][1..] == first[k + 1..];
    assert array_to_set(first[k..]) == {first[k]} + array_to_set(first[k + 1..]);
    firstSet := {first[k]} + firstSet;
    i := k;
    assert firstSet == array_to_set(first[i..]);
  }

  var secondSet: set<int> := {};
  var j := second.Length;
  while j > 0
    invariant 0 <= j <= second.Length
    invariant secondSet == array_to_set(second[j..])
    decreases j
  {
    var k := j - 1;
    assert 0 <= k < second.Length;
    assert secondSet == array_to_set(second[k + 1..]);
    assert |second[k..]| > 0;
    assert second[k..][0] == second[k];
    assert second[k..][1..] == second[k + 1..];
    assert array_to_set(second[k..]) == {second[k]} + array_to_set(second[k + 1..]);
    secondSet := {second[k]} + secondSet;
    j := k;
    assert secondSet == array_to_set(second[j..]);
  }

  assert first[0..] == first[..];
  assert second[0..] == second[..];
  assert firstSet == array_to_set(first[..]);
  assert secondSet == array_to_set(second[..]);
  are_same_set := firstSet == secondSet;
}
\end{lstlisting}
\section*{Final Program}
\begin{lstlisting}
function array_to_set(arr: seq<int>): set<int>
{
  if |arr| == 0 then {}
  else {arr[0]} + array_to_set(arr[1..])
}

method p_6_10_same_set_array(first: array<int>, second: array<int>) returns (are_same_set: bool)
	ensures are_same_set == (array_to_set(first[..]) == array_to_set(second[..]))
{
  var firstSet: set<int> := {};
  var i := first.Length;
  while i > 0
    invariant 0 <= i <= first.Length
    invariant firstSet == array_to_set(first[i..])
    decreases i
  {
    var k := i - 1;
    assert 0 <= k < first.Length;
    assert firstSet == array_to_set(first[k + 1..]);
    assert |first[k..]| > 0;
    assert first[k..][0] == first[k];
    assert first[k..][1..] == first[k + 1..];
    assert array_to_set(first[k..]) == {first[k]} + array_to_set(first[k + 1..]);
    firstSet := {first[k]} + firstSet;
    i := k;
    assert firstSet == array_to_set(first[i..]);
  }

  var secondSet: set<int> := {};
  var j := second.Length;
  while j > 0
    invariant 0 <= j <= second.Length
    invariant secondSet == array_to_set(second[j..])
    decreases j
  {
    var k := j - 1;
    assert 0 <= k < second.Length;
    assert secondSet == array_to_set(second[k + 1..]);
    assert |second[k..]| > 0;
    assert second[k..][0] == second[k];
    assert second[k..][1..] == second[k + 1..];
    assert array_to_set(second[k..]) == {second[k]} + array_to_set(second[k + 1..]);
    secondSet := {second[k]} + secondSet;
    j := k;
    assert secondSet == array_to_set(second[j..]);
  }

  assert first[0..] == first[..];
  assert second[0..] == second[..];
  assert firstSet == array_to_set(first[..]);
  assert secondSet == array_to_set(second[..]);
  are_same_set := firstSet == secondSet;
}
\end{lstlisting}
\section*{Total Token Usage}
\textbf{Input tokens: }434
\\\textbf{Output tokens: }6186
\\\textbf{Reasoning tokens: }5633
\\\textbf{Sum of `total tokens': }6620
\section*{Experiment Timings}
\textbf{Overall Experiment} started at 1780579889298, ended at 1780580070710, lasting 181412ms (181.41 seconds)
\\\textbf{Iteration \#1} started at 1780579889339, ended at 1780580070705, lasting 181366ms (181.37 seconds)
\end{document}
